\documentclass{article}
\usepackage{amsmath,amssymb,amsthm,xcolor}
\usepackage[margin=1in]{geometry}

\newtheorem{theorem}{Theorem}
\newtheorem{conjecture}{Conjecture}
\newtheorem{definition}{Definition}
\newtheorem{lemma}{Lemma}
\newtheorem{remark}{Remark}

\title{Cancelling Translated Copies of Points in Convex Position}
\date{}

\begin{document}

\maketitle

\section*{Problem Statement}

\begin{definition}
A finite set $S$ of points in the plane is in \emph{convex position} if every point of $S$ is a vertex of its convex hull.
\end{definition}

\begin{definition}[Setup]
Let $S$ be a finite set of $n \geq 1$ points in $\mathbb{R}^2$ in convex position.
Let $T$ be a finite multiset of translating vectors in $\mathbb{R}^2$, each assigned a sign
$\sigma(t) \in \{+1, -1\}$.

For every point $p \in \mathbb{R}^2$, define
\[
  \sigma(p) \;=\; \sum_{\substack{t \in T \\ \exists\, s \in S: \, p = t + s}} \sigma(t).
\]
\end{definition}

\begin{theorem}[Lower Bound -- Known]
Under the above setup, there are at least $n$ points $p \in \mathbb{R}^2$ for which $\sigma(p) \neq 0$.
\end{theorem}

\begin{conjecture}[Main -- To Be Proved]
If there are exactly $n$ points $p \in \mathbb{R}^2$ for which $\sigma(p) \neq 0$, then $S$ is the set of vertices of a parallelogram (in particular, $n = 4$).
\end{conjecture}

\section*{Goal}

Prove the conjecture from scratch.  That is, assuming the lower bound holds (or re-proving it as part of the argument), show that equality $|\{p : \sigma(p) \neq 0\}| = n$ forces $S$ to be a parallelogram (and hence $n = 4$).

\begin{remark}
The lower bound follows from convex geometry: the vertices of the convex hull of $S + T := \{t + s : t \in T, s \in S\}$ (counted with signed multiplicity) cannot all cancel.  The content of the conjecture is that the extremal case forces $S$ to be a parallelogram.
\end{remark}

\end{document}
